\documentclass[11pt]{article}
\usepackage[utf8]{inputenc}
\usepackage[T2A]{fontenc}
\usepackage[russian]{babel}
\usepackage{amsmath,amssymb}
\usepackage[margin=2.5cm]{geometry}
\usepackage{hyperref}
\usepackage{enumitem}

\title{Федеративное обучение: основы и применение к парку роботов}
\author{}
\date{11 сентября 2026}

\begin{document}
\maketitle

\section{Определение}

Федеративное обучение (FL, \emph{federated learning}) — парадигма обучения, в которой
модель обучается на данных, распределённых по множеству устройств (клиентов), без
централизованного сбора самих данных. Термин и первый практический алгоритм ввели McMahan с соавторами в работе
\emph{Communication-Efficient Learning of Deep Networks from Decentralized Data}
\cite{mcmahan2017}, где сформулированы три
определяющих свойства задачи: клиентские данные \textbf{non-IID} (распределены
неоднородно между клиентами), клиентов \textbf{много} (массовая распределённость), и
на каждом раунде обучения доступна лишь \textbf{часть} клиентов (частичное участие).

\section{Алгоритм FedAvg}

Обучение организовано раундами. На раунде $t$:
\begin{enumerate}
  \item Сервер рассылает текущие веса $w_t$ подвыборке клиентов.
  \item Каждый клиент $k$ выполняет несколько шагов SGD на своих локальных данных,
        получая $w_t^k$.
  \item Клиент отправляет серверу обновление — веса $w_t^k$ или дельту
        $\Delta_k = w_t^k - w_t$. \textbf{Сырые данные клиента никогда не покидают
        устройство.}
  \item Сервер агрегирует взвешенным средним по объёму локальных данных $n_k$:
        \[
          w_{t+1} = \sum_{k} \frac{n_k}{n} \, w_t^k, \qquad n = \sum_k n_k.
        \]
\end{enumerate}

Важно: сокрытие сырых данных — необходимое, но \emph{не достаточное} условие
приватности. По пересылаемым обновлениям $\Delta_k$ можно частично реконструировать
исходные данные (gradient inversion) — это мотивирует дополнительные слои защиты
(дифференциальная приватность, secure aggregation), рассмотренные в отдельных
разделах этого ресёрча.

\section{Проблема non-IID и способы её смягчения}

Ванильный FedAvg деградирует при сильной неоднородности клиентских распределений.
Основные направления коррекции:
\begin{itemize}
  \item момент-based коррекция локального SGD \cite{momentum2023};
  \item адаптивная агрегация через обучение с подкреплением \cite{feddrl2022};
  \item клиентская адаптация модели к локальному сдвигу распределения
        \cite{clientadapt2020}.
\end{itemize}
Общее ограничение этого семейства: каждый приём требует пересылки дополнительной
статистики о клиентском распределении, что расширяет поверхность потенциальной
утечки — прямое противоречие с целью приватности проекта.

Свежая работа \cite{fltaxonomy2025} даёт единую таксономию: топология
(централизованная / децентрализованная / иерархическая) и режим агрегации
(синхронный / асинхронный) — она используется как индекс в остальных разделах этого
ресёрча (архитектура флота, MPC, методы агрегации, топологии связи).

\section{Применение к парку VLA-роботов}

Каждый робот выступает клиентом FL. Локально он либо (а) дообучает VLA-политику на
своих демонстрациях, либо (б) считает локальную статистику калибровки
(non-conformity score для conformal prediction). В обоих случаях данные дома —
планировка, объекты, присутствие людей — никогда явно не покидают робота; покидают
только веса/дельты модели или скалярные статистики. Ключевая сложность именно этого
сценария — сильный non-IID между домами (кухня vs спальня, разные бытовые привычки),
что делает выбор между ванильным FedAvg и non-IID-коррекцией практическим, а не
теоретическим вопросом: коррекция даёт лучшую sample efficiency ценой пересылки
дополнительной статистики, которую нужно защищать отдельно.

\section{Соседние области}

\emph{Distributed SGD / parameter server} — исторический предшественник FedAvg,
перенос почти прямой; отличие FL — non-IID и partial participation как
первоклассные ограничения. \emph{Split learning} — альтернатива для случаев, где даже
обновления модели слишком чувствительны: модель физически разрезается между
устройством и сервером. Перенос дорогой (redesign архитектуры VLA), и разрезание
добавляет задержку в контур управления — риск для латентности реального времени.

\section{Тезаурус}

\begin{description}
  \item[Федеративное обучение (FL)] Обучение модели на данных, распределённых по
    множеству устройств, без централизованного сбора самих данных — на сервер едут
    только обновления модели.
  \item[Клиент] Устройство (в этом проекте — робот), локально обучающееся на своих
    данных и отправляющее обновления координатору.
  \item[Раунд] Один цикл: рассылка весов $\to$ локальное дообучение $\to$ отправка
    обновлений $\to$ агрегация на сервере.
  \item[FedAvg] Базовый алгоритм FL: сервер усредняет клиентские веса, взвешивая по
    объёму локальных данных $n_k$ каждого клиента.
  \item[Non-IID] Свойство клиентских данных быть неоднородно распределёнными между
    устройствами (например, разные комнаты дают разные распределения наблюдений) —
    основная причина деградации ванильного FedAvg.
  \item[Partial participation] Не все клиенты доступны на каждом раунде обучения —
    сервер работает с подвыборкой.
  \item[Gradient inversion] Атака, восстанавливающая исходные данные клиента по
    пересылаемому обновлению (градиенту/дельте весов) — причина, по которой сокрытия
    сырых данных недостаточно для приватности.
  \item[Дифференциальная приватность (DP)] Формальная гарантия, ограничивающая, сколько
    информации об отдельном примере можно извлечь из выхода алгоритма (добавлением
    контролируемого шума).
  \item[Split learning] Альтернатива FL, при которой сама модель физически разрезается
    между устройством и сервером, а не переобучается целиком на каждом клиенте.
\end{description}

\begin{thebibliography}{9}
\bibitem{mcmahan2017} B.~McMahan, E.~Moore, D.~Ramage, S.~Hampson, B.~Agüera y Arcas.
  \emph{Communication-Efficient Learning of Deep Networks from Decentralized Data}.
  AISTATS, 2017.
\bibitem{momentum2023} \emph{Momentum Benefits Non-IID Federated Learning Simply and
  Provably}. arXiv:2306.16504, 2023.
\bibitem{feddrl2022} \emph{FedDRL: Deep Reinforcement Learning-based Adaptive
  Aggregation for Non-IID Data in Federated Learning}. arXiv:2208.02442, 2022.
\bibitem{clientadapt2020} \emph{Client Adaptation improves Federated Learning with
  Simulated Non-IID Clients}. arXiv:2007.04806, 2020.
\bibitem{fltaxonomy2025} \emph{Federated Learning Survey: A Multi-Level Taxonomy of
  Aggregation Techniques, Experimental Insights, and Future Frontiers}.
  arXiv:2511.22616, 2025.
\end{thebibliography}

\end{document}
